% !TeX program = xelatex
\documentclass[11pt]{amsart}
\usepackage{fontspec}
\usepackage{polyglossia}
\setdefaultlanguage{russian}
\setotherlanguage{english}
\setmainfont{Libertinus Serif}
\setmonofont{DejaVu Sans Mono}
\usepackage{microtype}
\setlength{\emergencystretch}{2em}
\usepackage[margin=1in]{geometry}
\usepackage{amsmath,mathtools,unicode-math,tikz,listings,booktabs,tabularx,array,hyperref}
\setmathfont{Libertinus Math}
\usetikzlibrary{arrows.meta,positioning}
\hypersetup{hidelinks,pdftitle={Критерий Райснера}}
\newtheorem{lemma}{Лемма}[section]
\newtheorem{theorem}{Теорема}[section]
\lstdefinestyle{leanstyle}{basicstyle=\ttfamily\footnotesize,
  columns=fullflexible,keepspaces=true,breaklines=true,frame=none,
  showstringspaces=false,literate={σ}{{$\sigma$}}1 {ℤ}{{$\mathbb Z$}}1
    {ℕ}{{$\mathbb N$}}1 {↔}{{$\leftrightarrow$}}1 {∀}{{$\forall$}}1
    {∈}{{$\in$}}1 {→}{{$\to$}}1 {←}{{$\leftarrow$}}1
    {∧}{{$\land$}}1 {∅}{{$\varnothing$}}1}
\title{Обзор Lean-проекта. \\ \tiny{Критерий Райснера}}
\author{Лев Ерченков}
\date{}
\begin{document}
\begin{abstract}
В данном файле изложен обзор теоремы из комбинаторной алгебры, которую автор
формализовывал в Lean. Обзор состоит из формулировки утверждения и описания
Lean-формализации.
\end{abstract}
\maketitle

\section{Обзор утверждения: критерий Райснера}
Пусть $\mathcal{K}$ — симплициальный комплекс\footnote{Напомним, что это семейство
подмножеств множества вершин, замкнутое относительно перехода к
подмножествам; пустое множество также принадлежит комплексу.} на конечном
множестве вершин $V$, а $\mathbf{k}$ — поле. \textit{Кольцом
Стэнли--Райснера} комплекса $\mathcal{K}$ над $\mathbf{k}$ называется
фактор-кольцо
\[
  \mathbf{k}[\mathcal{K}]
  :=\mathbf{k}[x_v\mid v\in V]/I_{\mathcal{K}},\qquad
  I_{\mathcal{K}}:=\langle x_\sigma\mid \sigma\subseteq V,
      \ \sigma\notin\mathcal{K}\rangle,
  \qquad x_\sigma:=\prod_{v\in\sigma}x_v.
\]
Иными словами, $I_{\mathcal{K}}$ порождён мономами,
носители которых не являются гранями $\mathcal{K}$.

\textit{Глубиной} локального кольца $(R,\mathfrak n)$ называется максимальная длина
регулярной последовательности $a_1,\ldots,a_r$ элементов $\mathfrak n$;
это означает, что для каждого $j$ элемент $a_j$ не является делителем нуля
в фактор-кольце $R/(a_1,\ldots,a_{j-1})$.

Кольцо Стэнли--Райснера $\mathbf{k}[\mathcal{K}]$ называется
\textit{кольцом Коэн--Маколея}, если каждая его
локализация по простому идеалу является локальным кольцом с равными глубиной
и размерностью Крулля \cite[Definitions~10.103.1 and 10.104.6]{Stacks}.
Для грани $\sigma$ \textit{линк} определяется как
\[
  \operatorname{link}_{\mathcal{K}}(\sigma)
  =\{\tau\in\mathcal{K}: \tau\cap\sigma=\varnothing,
                    \ \tau\cup\sigma\in\mathcal{K}\}.
\]
Критерий Райснера утверждает \cite[Theorem~2.5]{Panov}:
\begin{theorem}
Следующие условия эквивалентны:
\begin{itemize}
  \item $\mathbf{k}[\mathcal{K}]$ — кольцо Коэн--Маколея;
  \item $\widetilde H_i(\operatorname{link}_{\mathcal{K}}(\sigma);\mathbf{k})=0$
    для каждой грани $\sigma\in\mathcal{K}$ и каждого $i\in\mathbb Z$,
    удовлетворяющего $i<\dim\operatorname{link}_{\mathcal{K}}(\sigma)$.
\end{itemize}
\end{theorem}
Условие проверяется для каждой грани, в том числе для пустой грани
$\varnothing$, чей линк равен самому $\mathcal{K}$. Используется
аугментированые гомологии: пустая грань лежит в степени $-1$, а группы
ниже $-1$ равны нулю.

В Lean условие Коэн--Маколея записано как \texttt{IsCohenMacaulay}, а
условие о занулении гомологии линков как
\texttt{LinkHomologyVanishesBelowDimension}. Ниже представлено объявление теоремы
из файла \texttt{Reisner/ReisnerCriterion.lean}; контекст переменных и
локальный экземпляр \texttt{IsNoetherianRing} заданы в этом файле. Объявление
компилируется, но доказательство пока не завершено: в нём стоит
\texttt{sorry}.

\begin{lstlisting}[style=leanstyle]
theorem reisner_criterion :
    IsCohenMacaulay (K.FaceRing k) ↔
      K.LinkHomologyVanishesBelowDimension k := by
  sorry
\end{lstlisting}

\section{Схема доказательства критерия}

Обозначим $A=\mathbf{k}[\mathcal{K}]$,
$\mathfrak m=(x_v:v\in V)$ и
$d=\max\{|\sigma|:\sigma\in\mathcal{K}\}=\dim A$.
Доказательство сводит Коэн-Маколеевость кольца к занулению
когомологий, а затем описываются эти когомологии через линки. Опишем план доказательства в обе стороны.

Для вычислений нам понадобится \textit{комплекс Чеха} на переменных $x_v$, $v\in V$, — это коцепной
комплекс с членами
\[
  C^q_{\check C}(A;x_v)
  =\prod_{\substack{F\subseteq V\\|F|=q}} A_{x_F},
  \qquad x_F=\prod_{v\in F}x_v,
  \qquad 0\leqslant q\leqslant |V|.
\]
Здесь $x_\varnothing=1$, поэтому $C^0_{\check C}(A;x_v)=A$;
дифференциалы являются знакопеременными суммами естественных отображений
локализации $A_{x_F}\to A_{x_G}$ при $F\subset G$. Так как $A$ нётерово,
когомологии этого комплекса совпадают с локальными когомологиями
$H^q_{\mathfrak m}(A)$.

\paragraph{\textbf{Коэн--Маколеевость $\Rightarrow$ условие на линки.}}
Если $A$ — кольцо Коэн--Маколея, то локальное кольцо $A_{\mathfrak m}$ имеет
глубину $d$. Критерий глубины через локальные когомологии даёт
$H^q_{\mathfrak m}(A)=0$ при $q<d$ \cite[Lemma~47.11.1]{StacksDepth}. Кольцо Стэнли-Райснера $\mathbb{Z}^n$-градуировано, применяем
формулу Хохстера для когомологий в степени $-\mathbf 1_\sigma$
:
\[
  \bigl(H^q_{\mathfrak m}(A)\bigr)_{-\mathbf 1_\sigma}
  \cong \widetilde H^{\,q-|\sigma|-1}
       (\operatorname{link}_{\mathcal{K}}(\sigma);\mathbf{k}).
\]
(См. \cite[Theorem~5.1]{Welker}.)
Здесь $\mathbf 1_\sigma$ — вектор, равный $1$ на вершинах грани $\sigma$
и $0$ на остальных вершинах. Для требуемой степени $i\geqslant -1$ берём
$q=|\sigma|+i+1$. Если
$i<\dim\operatorname{link}_{\mathcal{K}}(\sigma)$, то
$\dim\operatorname{link}_{\mathcal{K}}(\sigma)\leqslant d-|\sigma|-1$ (факт, верный для произвольного симплициального комплекса),
следовательно $q<d$ и когомологии линка в этой размерности зануляются. Поскольку коэффиценты беруется в поле $\mathbf{k}$, то гомологии тоже зануляются. Степени ниже $-1$ зануляются по принятому
аугментированному соглашению.

\paragraph{\textbf{Условие на линки $\Rightarrow$  Коэн--Маколеевость.}}
Сначала условие на зануление гомологий линков влечёт чистоту $\mathcal{K}$. Это доказывается
индукцией по размерности. Для чистого комплекса
\[
  \dim\operatorname{link}_{\mathcal{K}}(\sigma)=d-|\sigma|-1.
\]
Тогда условие на все линки равносильно занулению групп локальных когомологий
$H^q_{\mathfrak m}(A)$ при $q<d$. Критерий глубины показывает, что
локализация $A_{\mathfrak m}$ является кольцом Коэн--Маколея. Затем, с помощью критерия
для градуированных колец переносим этот результат на всё кольцо $A$
\cite[Item~2.1.27]{BrunsHerzog}.

Комбинаторный переход в этой схеме уже формализован через когомологии
фактического комплекса Чеха. Сравнение их с локальными когомологиями в
положительных степенях и остальные недоказанные переходы перечислены в конце.

Приведём небольшой пример. Возьмём $\mathcal{K}=\partial\Delta^2$: три вершины и три ребра, но без внутренней
двумерной грани. Геометрически это окружность, $\dim\mathcal{K}=1$. Кольцо
Стэнли--Райснера имеет вид
\[
  \mathbf{k}[\mathcal{K}]\cong
  \mathbf{k}[x_1,x_2,x_3]/(x_1x_2x_3).
\]
\begin{center}
\small
\begin{tabularx}{\textwidth}{@{}p{2cm}p{3.5cm}p{2.6cm}X@{}}
\toprule
\textbf{Грань $\sigma$} & \textbf{Линк} & \textbf{Размерность} &
\textbf{Требуемое зануление групп гомологии} \\
\midrule
Пустое $\varnothing$ & $\mathcal{K}$ (окружность) & $1$ &
Степени $-1,0$; обе группы равны нулю. Верхняя группа
$\widetilde H_1(\mathcal{K};\mathbf{k})\cong\mathbf{k}$ разрешена. \\
Вершина $\{v_1\}$ & Две изолированные вершины & $0$ &
Только степень $-1$; группа равна нулю. Группа
$\widetilde H_0\cong\mathbf{k}$ разрешена. \\
Ребро $\{v_1,v_2\}$ & $\{\varnothing\}$ & $-1$ &
Степени $i<-1$ равны нулю. \\
\bottomrule
\end{tabularx}
\end{center}
Линки остальных вершин и рёбер имеют тот же вид. Значит, условие выполнено
для каждой грани, и критерий заключает, что $\mathbf{k}[\mathcal{K}]$ является
кольцом Коэн--Маколея.

\section{Промежуточная лемма и её Lean-формализация}

Следующая лемма доказана в Lean. Положим
$d(\mathcal{K})=\max\{|\sigma|:\sigma\in\mathcal{K}\}$ — максимальное
число вершин в грани. Через
$C^\bullet_{\check C}(\mathbf{k}[\mathcal{K}])$ обозначим
комплекс Чеха кольца
Стэнли--Райснера, реализованный в проекте как
\texttt{K.cechComplex k}.

\begin{lemma}[Зануление комплекса Чеха и гомологии линков]\label{lem:cech-link-homology}
Пусть $\mathcal{K}$ содержит пустую грань и задан на конечном
линейно упорядоченном множестве вершин $V$, а $\mathbf{k}$ — поле.
Следующие условия Эквивалентны:
\begin{enumerate}
\item $H^q(C^\bullet_{\check C}(\mathbf{k}[\mathcal{K}]))=0$ для всех
      $q\in\mathbb N$ с $q<d(\mathcal{K})$;
\item $\mathcal{K}$ чистый и для каждой грани $\sigma\in\mathcal{K}$ группы
      редуцированных гомологий её линка зануляются ниже размерности линка:
      \[
        \forall i\in\mathbb Z,\qquad
        i<\dim\operatorname{link}_{\mathcal{K}}(\sigma)
        \ \Longrightarrow\
        \widetilde H_i(\operatorname{link}_{\mathcal{K}}(\sigma);\mathbf{k})=0.
      \]
\end{enumerate}
Здесь чистота означает, что все максимальные грани имеют $d(\mathcal{K})$ вершин;
\end{lemma}

Критерий Райснера формулирует комбинаторное условие через группы гомологий
каждого линка. Эта лемма переводит его (вместе с чистотой) в единое условие на
низкие степени комплекса Чеха. Так комбинаторную часть доказательства можно
свести к занулению комплекса Чеха и затем сопоставить его с локальными
когомологиями и глубиной кольца. Сама лемма такого сопоставления не доказывает:
для полного критерия ещё нужны вывод чистоты из условия на все линки и переход
от алгебраического свойства Коэн--Маколея к нужному занулению комплекса Чеха
или к занулению локальных когомологий. Итоговая теорема
\texttt{reisner\_criterion} в Lean пока содержит \texttt{sorry}.

Ниже приведено точное объявление с доказательством. Заметим, что в Lean условие нулевой
группы гомологии выражено через \texttt{Subsingleton}.
\begin{lstlisting}[style=leanstyle]
namespace Reisner.Complex
universe u v

variable {V : Type u} (K : Reisner.Complex V)
  (k : Type v) [Field k] [Fintype V] [LinearOrder V]

theorem cechCohomology_vanishes_iff_pure_and_linkHomology_vanishes :
    (∀ q : ℕ, q < K.maxFaceCard →
      Subsingleton ((K.cechComplex k).homology q)) ↔
    (∀ F : K.facets, F.1.card = K.maxFaceCard) ∧
      (∀ (σ : Finset V) (hσ : σ ∈ K.faces) (i : ℤ),
        i < (K.link σ hσ).dimension →
          Subsingleton (((K.link σ hσ).simplicialChainComplexInt k).homology i)) := by
  rw [K.cechCohomology_vanishes_iff_pure_and_linkCohomology_vanishes k]
  apply and_congr_right
  intro _
  refine forall_congr' fun σ => forall_congr' fun hσ => forall_congr' fun i =>
    imp_congr_right fun _ => ?_
  rw [← ModuleCat.isZero_iff_subsingleton,
    ← ModuleCat.isZero_iff_subsingleton]
  exact (K.link σ hσ).simplicialCochainInt_homology_isZero_iff_chainInt_homology_isZero k i

end Reisner.Complex
\end{lstlisting}

Доказательство сначала переписывает зануление комплекса Чеха через уже
установленное зануление когомологий линков. Затем с помощью двойственности
аугментированных цепных и коцепных комплексов заменяются когомологии линков на
гомологии, сохраняя целую степень, включая $-1$.
\section{Что сделано в Lean}

В Lean формализовано кольцо Стэнли--Райснера и использовано алгебраическое
свойство Коэн--Маколея. Точная формулировка критерия записана для конечного
множества вершин и поля; объявление компилируется, но доказательство пока
содержит \texttt{sorry}.

Уже доказано, что зануление когомологий фактического комплекса Чеха ниже
$d(\mathcal{K})$ эквивалентно чистоте $\mathcal{K}$ и занулению
гомологий каждого линка ниже его собственной размерности. Также доказана
двойственность аугментированных цепных и коцепных комплексов во всех целых
степенях, включая степень $-1$; степени ниже $-1$ зануляются по определению.


\begin{thebibliography}{9}
\bibitem{Panov}
T. Panov, \emph{Cohomology of face rings, and torus actions},
arXiv:math/0506526v3 (2007), Sections~2.1--2.2, Theorem~2.5.
\url{https://arxiv.org/html/math/0506526v3}.

\bibitem{Stacks}
The Stacks Project, \emph{Algebra}, Definitions~10.103.1 and~10.104.6,
Tags 00N3 and 00NC.
\url{https://stacks.math.columbia.edu/tag/00N3};
\url{https://stacks.math.columbia.edu/tag/00NC}.

\bibitem{StacksDepth}
The Stacks Project, Lemma~47.11.1, depth and the first nonzero local
cohomology group, Tag 0AVZ.
\url{https://stacks.math.columbia.edu/tag/0AVZ}.

\bibitem{Welker}
V. Welker, \emph{Which properties of Stanley--Reisner rings and simplicial
complexes are topological?}, Section~5, Theorem~5.1.
\url{https://www.mathematik.uni-marburg.de/~welker/preprints/top_survey.pdf}.

\bibitem{BrunsHerzog}
W. Bruns and J. Herzog, \emph{Cohen--Macaulay Rings}, revised edition,
Cambridge Studies in Advanced Mathematics~39, Cambridge University Press,
1998, Item~2.1.27.
\end{thebibliography}

\end{document}
